\section{A Gram matrix for the seed constants}\label{sec:gram}

We construct $B$ candidate multipliers. Each produces a vector of
$B$ seed constants, one for each odd grid. These vectors form
the columns of a rational matrix $\mathcal M$. If $\mathcal M$
is nonsingular, a rational combination of its columns gives
$e_B$; clearing denominators gives the seed vector
\eqref{eq:seed-target} with integer coefficients on the multipliers.
We adapt the Gram representation of~\cite[Section~3]{Tcount},
first showing the choice with two grids.

\subsection{A worked example with two grids}\label{sec:worked-gram}

Take $B=2$, $m=1$, $r=2$ and $a=70$, the parameters that will
give two irrational values through $\beta(70)$. There is no
differentiation, since $k=0$. The grids have $b=1,3$.
We want $q_1(1)=0$ and $q_3(1)=\eta\ne0$. Indeed, for any
possible irrational value $\beta(\nu)$, the fixed combination
$3^\nu F_{n,1}-F_{n,3}$ cancels its coefficient at every $n$.
Its seed constant is then $-\eta\ne0$.

At $n=1$, put $z=t+1$. The two kernels are
\[
 H_1(t)=\frac1z,\qquad
 H_3(t)=\frac1{z-2/3}-\frac1z+\frac1{z+2/3}.
\]
Expansion at $t=0$ gives $q_b(1)=\Res_{t=0}T_1(t)H_b(t)$.
We will choose independent polynomials whose pairwise products,
integrated against one positive weight, give the seed matrix
up to a common sign. The line $z=iy$ lies midway between
the poles $t=0,-2$ and turns every symmetric factor
$z^2-\alpha^2$ into $-(y^2+\alpha^2)$, with fixed sign.
Write $x=y^2$. On this line,
$H_b(-1+iy)=iy\,g_b(x)$, with
\[
 g_1(x)=-\frac1x,\qquad
 g_3(x)=\frac1x-\frac2{x+4/9}.
\]
Multiplying by their common denominator $x(9x+4)$ gives
two independent polynomials,
\[
 h_1(x)=-9x-4,\qquad h_3(x)=4-9x.
\]
To obtain products $h_bh_j$ in the integral, choose multipliers
whose values on this line reproduce the polynomials:
$W_{1,j}(-1+iy)=h_j(y^2)$. Substituting $-(t+1)^2$ for $x$
and homogenizing in $n$ gives
\[
 W_{n,1}(t)=9(t+n)^2-4n^2,\qquad
 W_{n,3}(t)=9(t+n)^2+4n^2.
\]
These are integer polynomials, even about $-n$ and homogeneous
of degree $2$.

Let $\mathcal M$ have rows indexed by the seed constants and
columns indexed by these multipliers. Closing the midpoint
contour clockwise gives $\mathcal M=-\mathcal N$, where
\[
 \mathcal N_{b,j}=\int_{-\infty}^{\infty}
              w(y)h_b(y^2)h_j(y^2)\,dy,\qquad b,j\in\{1,3\},
\]
and the positive weight is
\[
 w(y)=\frac{2^{85}3^{16}}{18\pi}
 \frac{\displaystyle
  \prod_{\alpha\in\{4/3,\,2,\,8/3,\,10/3,\,4,\,14/3\}}
                         (y^2+\alpha^2)}
      {(1+y^2)^{70}}.
\]
The numerator's central zero cancels the kernel pole at $z=0$;
its factor $z^2-4/9$ cancels the paired poles at $z=\pm2/3$.
On the midpoint line, $P_1$ is $-iy$ times a positive product,
since it has seven symmetric root pairs. Multiplying by the
kernel's $iy$ factor removes that sign, and the clockwise
contour supplies the minus sign in $\mathcal M=-\mathcal N$.
Section~\ref{sec:positive-weight} gives the general calculation.

Put $\mu_\ell=\int_{-\infty}^{\infty}w(y)y^{2\ell}\,dy$ for
$\ell=0,1,2$. The matrix is explicitly
\[
 \mathcal N=
 \begin{pmatrix}
 16\mu_0+72\mu_1+81\mu_2&-16\mu_0+81\mu_2\\
 -16\mu_0+81\mu_2&16\mu_0-72\mu_1+81\mu_2
 \end{pmatrix}.
\]
Its determinant is
\[
 \det\mathcal N=5184(\mu_0\mu_2-\mu_1^2)>0.
\]
The inequality is strict Cauchy--Schwarz: $y^2$ is not constant
under the positive weight $w$. The rational combination
\[
 W_n=\mathcal N_{1,3}W_{n,1}-\mathcal N_{1,1}W_{n,3}
\]
has seed vector $(0,\det\mathcal N)^{\mathsf T}$. Its coefficients
are rational because $\mathcal M$ is a matrix of rational residues.
Clearing their fixed denominators gives an integer multiplier
with seed vector $(0,\eta)$, $\eta\ne0$. We now make the same
choice with $B$ independent polynomials.

\subsection{Moving the derivatives onto the kernels}

At $n=1$, the poles are $0$ and $-2$. For a proper rational
function $V$ with these poles and antisymmetric about $-1$, its partial-fraction
coefficients satisfy $A_{\nu,1}=(-1)^{\nu+1}A_{\nu,0}$.
For odd $b\le M$, define
\[
 H_b(t)=\sum_{\substack{1\le j<2b\\j\text{ odd}}}
                           \frac{\ch(j)}{t+j/b}.
\]
Expansion at $0$ and \eqref{eq:q} give
$q_b=\Res_{t=0}V(t)H_b(t)$. If $V=T^{(k)}/k!$, repeated
integration by parts for residues gives
\begin{equation}\label{eq:residue}
 q_b=\Res_{t=0}T(t)\frac{H_b^{(k)}(t)}{k!}.
\end{equation}
The sign is positive because $k$ is even. Transferring the
derivatives to the kernel leaves the numerator of $T$ intact.

Put $c=(m+1)/2$ and $z=t+1$. The central kernel term is
$\ch(b)/z$. The terms at distances $\pm\alpha$,
$\alpha=2j/b$, have equal weights $\ch(b-2j)=\ch(b+2j)$.
On $z=iy$ we therefore have
\[
 \frac{H_b^{(k)}(-1+iy)}{k!}=iy\,g_b(y^2),
\]
where
\begin{equation}\label{eq:g}
 \begin{split}
 g_b(x)={}&\frac{(-1)^c\ch(b)}{x^c}\\
 &-2\sum_{j=1}^{(b-1)/2}\ch(b-2j)
 \frac{\displaystyle\sum_{\ell=0}^{c-1}(-1)^\ell
       \binom m{2\ell+1}\alpha^{m-2\ell-1}x^\ell}
      {(x+\alpha^2)^m},\qquad \alpha=\frac{2j}{b}.
 \end{split}
\end{equation}
This follows by combining $(iy-\alpha)^{-m}+(iy+\alpha)^{-m}$.
For $b=1$ the sum is empty: $g_1=-1/x$ when $m=1$, and
$g_1=1/x^2$ when $m=3$.

\subsection{Independent integer polynomials}

Let $\mathcal A$ be the distinct positive distances in
\eqref{eq:g}. In reduced form they are $v/d$ with $d\le M$
odd, $v$ even and $0<v<d$. Thus
\[
 N=|\mathcal A|=\frac{D-1}{2}.
\]
Use reduced numerators and denominators to define
\begin{equation}\label{eq:Q-h}
 Q(x)=\prod_{v/d\in\mathcal A}(d^2x+v^2),\qquad
 h_b(x)=x^cQ(x)^m g_b(x),\qquad J=mN.
\end{equation}
Each $h_b$ is an integer polynomial of degree at most $J$.
The central term contributes $(-1)^c\ch(b)Q^m$. For a paired
term with $\alpha=v/d$, clearing $(x+\alpha^2)^m$ leaves
the integer numerator coefficients
\[
 -2\ch(b-2j)(-1)^\ell\binom m{2\ell+1}
                  v^{m-2\ell-1}d^{m+2\ell+1}.
\]
After multiplication by $x^c$ and the other factors of $Q^m$,
its degree is at most $c+(c-1)+m(N-1)=mN$.

\begin{lemma}\label{lem:independence}
The polynomials $h_1,h_3,\ldots,h_M$ are linearly independent
over $\R$.
\end{lemma}
\begin{proof}
In a nonzero relation among the $g_b$, choose the largest $b$
with nonzero coefficient. If $b>1$, its pole at
$x=-((b-1)/b)^2$ has exact order $m$ and is absent from every
earlier kernel: it is farther from $0$ than any of their
negative poles. The order is exact because the corresponding
term of $H_b^{(k)}/k!$ has nonzero principal coefficient, and
$x=-z^2$ has nonzero derivative there. If $b=1$, then $g_1$
is nonzero. Thus no relation is possible. Multiplication by
the common denominator preserves independence.
\end{proof}

Homogenize the polynomials to obtain
\begin{align}
 W_{n,b}(t)
 &=n^{2J}h_b\!\left(-\frac{(t+n)^2}{n^2}\right)\notag\\
 &=\sum_{j=0}^J h_{b,j}(-1)^j n^{2J-2j}(t+n)^{2j}.
 \label{eq:W}
\end{align}
They have integer coefficients, are even about $-n$, and
have total degree $E=2J=m(D-1)$, as required in \eqref{eq:E}.
In particular,
\begin{equation}\label{eq:homogeneity}
 W_{p,b}(pt)=p^E W_{1,b}(t).
\end{equation}

\subsection{The common positive weight}\label{sec:positive-weight}

At $n=1$, the symmetric simple roots give
\[
 P_1(t)=z\prod_{\alpha\in\Lambda}(z^2-\alpha^2),
 \qquad z=t+1,\qquad |\Lambda|=\frac{DL-1}{2},
 \qquad\mathcal A\subseteq\Lambda.
\]
All $\alpha$ in $\Lambda$ are positive. Let $\mathcal M_{b,j}$
be the $b$th seed constant associated with $W_{1,j}$.
Every pole of $H_b^{(k)}/k!$ in $(-2,0)$, including the central
pole, has order at most $m$ and is canceled by $P_1^m$.
Thus the product in \eqref{eq:residue} has poles only at $0,-2$.

Close the upward line $t=-1+iy$ by a large right semicircle,
which encloses $0$ clockwise. Since $H_b^{(k)}(t)=O(t^{-m})$,
the integrand has degree at most
$mDL+E-2a-m<-1$ at infinity. The arc tends to zero.
On the midpoint line, $P_1^m$ contributes
$i^m(-1)^{|\Lambda|}y^m$ times a positive product.
The kernel contributes
$iy\,h_b(y^2)/(y^{2c}Q(y^2)^m)$.
Since $2c=m+1$, the powers of $y$ cancel. The clockwise
orientation supplies the final minus sign, giving
\begin{equation}\label{eq:Gram}
 \mathcal M_{b,j}=\frac{\sigma\kappa_1}{2\pi}
       \int_{-\infty}^{\infty}\omega(y)h_b(y^2)h_j(y^2)\,dy,
 \qquad \sigma=(-1)^{c+(DL-1)/2+1},
\end{equation}
where
\[
 \omega(y)=\frac{\displaystyle
                   \prod_{\alpha\in\Lambda}(y^2+\alpha^2)^m}
                  {(y^2+1)^a Q(y^2)^m}.
\]
This weight is strictly positive for all real $y$.
The central singularity has canceled, and the degree estimate
proves absolute convergence. Lemma~\ref{lem:independence}
therefore makes $\sigma\mathcal M$ positive definite: its
quadratic form is the integral of the square of a nonzero
polynomial against this weight.

The entries of $\mathcal M$ are rational by \eqref{eq:residue}.
Solve $\mathcal M\mathbf c=e_B$ over $\Q$ and clear denominators.
This gives integer coefficients $c_j$ and a nonzero rational
$\eta$ such that $\mathcal M\mathbf c=\eta e_B$.
Set
\[
 W_n=\sum_{\substack{1\le j\le M\\j\text{ odd}}}c_j W_{n,j}.
\]
This proves \eqref{eq:seed-target}. The coefficients remain fixed
as $n$ varies, so their sizes do not affect an exponential rate.
